% DI Atlas paper — v1 draft (Abstract + Introduction)
% Resource-paper framing: the atlas is the contribution; nulls characterize the resource.

\documentclass[11pt]{article}

\usepackage{amsmath,amssymb}
\usepackage{booktabs}
\usepackage{graphicx}
\usepackage{hyperref}
\usepackage[margin=1in]{geometry}
\usepackage{natbib}
\usepackage{lineno}
\usepackage{xcolor}

\title{A Pre-Registered Direction Instability Atlas\\
Across Three Perturbation Modalities}

\author{Elliot Tower\\
\texttt{elliot@elliottower.ai}}

\date{July 2026}

\linenumbers

\begin{document}
\maketitle

\begin{abstract}
When a drug or gene knockout is applied across many cell lines, the
resulting expression or morphological signatures can point in different
directions---a property quantified by cosine-based direction instability
(DI). A companion study established that DI is immune to the magnitude
confound that dominates alternative geometric metrics.  Here we construct a cross-modal DI atlas spanning three perturbation
modalities: LINCS L1000 transcriptomics (8{,}949 drugs), Tahoe-100M
single-cell transcriptomics (379 drugs), and JUMP Cell Painting
morphology (7{,}948 CRISPR knockouts, 12{,}590 ORF overexpressions).
We test nine pre-registered hypotheses about what DI captures. One passes: DI is
concordant across expression and morphology modalities (partial
$\rho = 0.19$, $n = 378$ compounds, $p = 2.1 \times 10^{-4}$).  The
remaining eight fail, and their failure pattern is informative.
Essential genes produce \emph{lower} DI ($\rho = -0.33$), inverting
the pre-registered prediction and revealing that functional constraint
suppresses context-dependence.  Knockout and overexpression DI for the
same gene are uncorrelated ($\rho = 0.01$, $n = 5{,}220$),
establishing DI as a property of the perturbation mechanism rather than
the target gene.  Expression variance, paralog count, and paralog
expression breadth each explain less than 0.4\% of DI variance after
confound control, ruling out three plausible proxy explanations.  The
atlas, pre-registration documents, and all analysis code are publicly
available.
\end{abstract}


%% ====================================================================
\section{Introduction}
%% ====================================================================

High-throughput perturbation screens now profile thousands of drugs and
genetic perturbations across panels of cell lines, producing
multi-dimensional signatures in transcriptomic
\citep{subramanian2017next}, morphological \citep{bray2016cell}, and
single-cell modalities \citep{tahoe2025}.  A central question in
perturbation biology is whether a compound or knockout acts the same way
across cellular contexts or produces context-dependent effects.  This
question has practical consequences: drugs with context-invariant
mechanisms are more likely to transfer from the cell line where they were
discovered to the patient tissue where they must work
\citep{tower2026confound}.

Direction instability (DI) quantifies this context-dependence as
$\mathrm{DI} = 1 - \overline{\cos(\mathbf{s}_i, \mathbf{s}_j)}$, the
mean pairwise cosine dissimilarity among a perturbation's signatures
across contexts.  Because cosine normalizes each signature to unit length before
comparison, DI is immune by construction to the magnitude
confound---systematic covariation between effect size and metric value.
Alternative geometric measures such as Grassmannian geodesic distance,
profile variability, and optimal transport cost are dominated by this
artifact \citep{tower2026confound}.  That immunity was the subject
of a companion study, which demonstrated magnitude confounding across
three perturbation modalities and showed that DI alone predicts
cross-cell-line mechanism transport (leave-one-out AUROC $= 0.986$ on
8{,}949 LINCS drugs).  The present study takes DI's validity as
established and asks a different question: what does DI capture, and
what does it fail to capture, across modalities and perturbation types?

We constructed a DI atlas spanning three measurement technologies and
four perturbation classes: drug perturbations measured by L1000
transcriptomics (LINCS, 8{,}949 compounds) and single-cell
transcriptomics (Tahoe-100M, 379 compounds); genetic knockouts measured
by Cell Painting morphology (JUMP-CP CRISPR, 7{,}948 genes); and
genetic overexpressions measured by Cell Painting morphology (JUMP-CP
ORF, 12{,}590 genes).  For each perturbation, the atlas provides
magnitude-corrected DI with bootstrap confidence intervals, enabling
downstream analyses that require a calibrated measure of
context-dependence.

To characterize this resource, we pre-registered nine hypotheses across
two batches before computing DI on new datasets (Batch~1 SHA
\texttt{0d07a01}; Batch~2A SHA \texttt{a17c125}).  The hypotheses test
whether DI is concordant across modalities (A1, A2), whether it
correlates with gene essentiality (C), whether it is an intrinsic
property of genes (D), and whether expression variance (E2), drug
selectivity (E3), or paralog buffering (E5a, E5b) can explain it.  Of these, one
passes its pre-registered criterion and eight fail.

The failures are the core result.  Cross-modal concordance (A2) confirms
that DI reflects biology rather than measurement artifact: compounds
whose transcriptomic effects vary across cell lines also produce
variable morphological profiles.  Every other test delimits what DI
measures.  Essential genes show \emph{low} DI, because functional
constraint produces stereotyped disruption regardless of cellular
identity.  Knockout and overexpression DI are unrelated, meaning DI
characterizes how a specific perturbation mechanism interacts with
context rather than indexing a stable property of the target gene.
Expression variance, paralog count, and paralog expression breadth
produce correlations too small to matter biologically (partial
$\rho < 0.07$ after confound control), ruling them out as explanations
for DI variation.

Together, these results establish DI as a distinct axis of perturbation
characterization.  It captures a specific interaction between
perturbation mechanism and cellular compensatory network that survives
cross-modal validation, resists reduction to existing gene annotations,
and is mechanism-specific rather than target-intrinsic.  The atlas
itself---DI values with confidence intervals across four perturbation
classes---is the primary contribution, available for integration with
downstream analyses such as drug mechanism prediction, gene function
annotation, and perturbation screen quality control.

The remainder of the paper is organized as follows.
Section~\ref{sec:atlas} describes the atlas as a resource: datasets,
DI computation, and magnitude correction.
Section~\ref{sec:concordance} presents the construct validity results
(A1 and A2), with PRISM viability as independent corroboration.
Section~\ref{sec:whatnot} reports the informative null and sign-flipped
results that delimit what DI measures (C, D, E2, E3, E5).
Section~\ref{sec:integrity} describes the pre-registration protocol and
deviations.
Section~\ref{sec:discussion} synthesizes the failure pattern into a
characterization of DI as an irreducible perturbation property.


%% ====================================================================
%% Placeholder sections (to be drafted)
%% ====================================================================

\section{The DI atlas as a resource}
\label{sec:atlas}

\textcolor{red}{[TO BE DRAFTED: datasets table, DI computation, magnitude correction via OLS residualization, bootstrap CIs, data availability.]}

\section{Construct validity: cross-modal concordance}
\label{sec:concordance}

\textcolor{red}{[TO BE DRAFTED: A2 as headline (rho=0.19, n=378, PASS). A1 as underpowered but effect-present (rho=0.25, n=76, FAIL CI gate). PRISM viability CV as exploratory corroboration (rho=0.36, n=75). A1/A2 asymmetry sentence. Do not oversell A2.]}

\section{What DI is not: informative nulls and the essentiality sign-flip}
\label{sec:whatnot}

\textcolor{red}{[TO BE DRAFTED: C sign-flip (post-hoc interpretation, label it). D null (mechanism-specific, not gene-intrinsic). E2 (real but negligible). E5a (confound removal textbook). E5b (tiny, gated). E3 (underpowered, magnitude confound confirmed). Frame as scope-delimitation of the atlas resource.]}

\section{Pre-registration integrity}
\label{sec:integrity}

\textcolor{red}{[TO BE DRAFTED: Two-batch protocol. SHA freezes. 7 deviations disclosed. DepMap version reconciliation (single paragraph). A1 matching correction (moved toward prereg, not away). CI gate rationale.]}

\section{Discussion}
\label{sec:discussion}

\textcolor{red}{[TO BE DRAFTED: Synthesis from V2 writeup. DI as irreducible axis. Resource-paper framing: the nulls characterize the resource. Limitations (PRISM exploratory, A1 underpowered, context definitions vary). Future work (Batch 2B / Shesha, GO enrichment, larger Tahoe overlap).]}


%% ====================================================================
%% References (to be replaced with .bib)
%% ====================================================================

\begin{thebibliography}{99}

\bibitem[Bray et~al., 2016]{bray2016cell}
Bray, M.-A., et~al. (2016).
Cell Painting, a high-content image-based assay for morphological profiling using multiplexed fluorescent dyes.
\textit{Nature Protocols}, 11:1757--1774.

\bibitem[Chandrasekaran et~al., 2024]{chandrasekaran2024cpjump1}
Chandrasekaran, S.N., et~al. (2024).
Three million images and morphological profiles of cells treated with matched chemical and genetic perturbations.
\textit{Nature Methods}, 21:1114--1121.

\bibitem[Chandrasekaran et~al., 2025]{chandrasekaran2025morphmap}
Chandrasekaran, S.N., et~al. (2025).
Morphological map of under- and overexpression of genes in human cells.
\textit{Nature Methods}.

\bibitem[Corsello et~al., 2020]{corsello2020prism}
Corsello, S.M., et~al. (2020).
Discovering the anticancer potential of non-oncology drugs by systematic viability profiling.
\textit{Nature Cancer}, 1:235--248.

\bibitem[Funk et~al., 2022]{funk2022phenotypic}
Funk, L., et~al. (2022).
The phenotypic landscape of essential human genes.
\textit{Cell}, 185(24):4634--4653.

\bibitem[Haghighi et~al., 2022]{haghighi2022crossmodal}
Haghighi, M., et~al. (2022).
High-dimensional gene expression and morphology profiles of cells across 28,000 genetic and chemical perturbations.
\textit{Nature Methods}, 19:1550--1557.

\bibitem[Subramanian et~al., 2017]{subramanian2017next}
Subramanian, A., et~al. (2017).
A next generation connectivity map: L1000 platform and the first 1,000,000 profiles.
\textit{Cell}, 171(6):1437--1452.

\bibitem[Tahoe Consortium, 2025]{tahoe2025}
Tahoe Therapeutics (2025).
Tahoe-100M: A giga-scale single-cell perturbation atlas for context-dependent gene function and cellular modeling.
\textit{bioRxiv}, 2025.02.20.639398.

\bibitem[Tower, 2026]{tower2026confound}
Tower, E. (2026).
Direction instability identifies drug mechanism transport while geometric alternatives measure magnitude.
\textit{In preparation}.

\end{thebibliography}

\end{document}
